% !PRINTED 388
\weberpage{388}

%<<CONT
it, as a normal divisor of $P$, also contains, when $(4,5,6)$ is a three-term cycle from $S'$ [by \S.\,102, (8), (10)]:
\[
\Pi_{1,2,3,5}\,(5,4,6)\,\Pi_{1,2,3,5} = (5,4,6)\,\Pi_{1,2,3,4}\,\Pi_{1,2,3,5}
\]
\[
= (5,4,6)(4,5)\,\Pi_{1,2,3,4},
\]
and consequently also
\[
S'(4,5)\,\Pi_{1,2,3,4} = S''\,\Pi_{1,2,3,4}.
\]
Now if
\[
\sigma = \begin{pmatrix} 1, & 2, & 3, & 4, & 5, & 6, & 7, & 8 \\ \alpha_1, & \alpha_2, & \alpha_3, & \alpha_4, & \alpha_5, & \alpha_6, & \alpha_7, & \alpha_8 \end{pmatrix}
\]
is a permutation from $S$ in which $\alpha_1, \alpha_2, \alpha_3, \alpha_4$ are arbitrarily given, then one can still choose the arrangement of the $\alpha_5, \alpha_6, \alpha_7, \alpha_8$ so that $\sigma$ belongs at will to $S'$ or to $S''$. If one chooses $\sigma$ in $S'$ and observes that then $S'' = S''\sigma^{-1}$, it follows that
\[
S''\,\sigma^{-1}\,\Pi_{1,2,3,4}\,\sigma = S''\,\Pi_{\alpha_1,\alpha_2,\alpha_3,\alpha_4}
\]
must also be contained in $Q$. But now we have also
\[
\Pi_{1,2,3,4}\,(3,5)(4,6)\,\Pi_{1,2,3,4} = (3,5)(4,6)\,\Pi_{1,2,5,6}\,\Pi_{1,2,3,4}
\]
\[
= (3,5)(4,6)(1,2)(3,4)(5,6)(7,8)\,\Pi_{1,2,7,8}
\]
\[
= (1,2)(3,6)(4,5)(7,8)\,\Pi_{1,2,7,8},
\]
and since $(1,2)(3,6)(4,5)(7,8)$ belongs to $S'$, $Q$ also contains $S'\,\Pi_{1,2,7,8}$, and consequently, as above, all $S'\,\Pi_{\alpha_1,\alpha_2,\alpha_3,\alpha_4}$, hence also all $S\,\Pi_{\alpha_1\alpha_2\alpha_3\alpha_4}$, and therefore also
\[
S\,\Pi_{\alpha_1,\alpha_2,\alpha_3,\alpha_4}\,\Pi_{\alpha_1,\alpha_2,\alpha_3,\alpha_4} = S,
\]
that is, $Q$ comprises the whole group $P$. From this it follows, in combination with 1.:

2.\ \gesperrt{If a normal divisor $Q$ of $P$ has, besides the identical permutation, any permutation whatever in common with $S$, then $Q$ is identical with $P$.}

It may now further be asked whether $Q$ can contain a permutation from one of the subsidiary groups of (1), that is, an element of the form $\sigma\,\Pi_{1,2,3,4}$, without being identical with $P$. If first $\sigma = 1$, so that $Q$ contains the element $\Pi_{1,2,3,4}$, then as a normal divisor of $P$ it also contains all the other $\Pi_{\alpha_1\alpha_2\alpha_3\alpha_4}$, as may be seen from the formula \S.\,102, (8), and with them also $\Pi_{1,2,3,4}\,\Pi_{1,2,3,5}\,\Pi_{1,2,3,4} = (4,5)$, and is therefore by 2. identical with $P$.
